\documentclass{article}
\usepackage[utf8]{inputenc}
\usepackage[T1]{fontenc}
\usepackage{geometry}
\usepackage{graphicx}
\usepackage{float}
\usepackage{hyperref}
\usepackage[english,indonesian]{babel}

\title{Optimasi Pemrosesan Template LaTeX Menggunakan Pendekatan Berbasis AI}
\author{Academic Support Agent}
\date{\today}

\begin{document}
\maketitle

\begin{abstract}

\end{abstract}

\chapter{PENDAHULUAN}


\chapter{TINJAUAN PUSTAKA}


% Original content (commented out):
% %%%%%%%%%%%%%%%%%%%%%%%%%%%%%%%%%%%%%%%%%%%%%%
%%                                          %%
%% USE THIS FILE TO SUBMIT YOUR SOLUTIONS   %%
%%                                          %%
%% You must have the usamts.tex file in     %%
%% the same directory as this file.         %%
%% You do NOT need to submit this file or   %%
%% usamts.tex with your solutions.  You     %%
%% only need to submit the output PDF file. %%
%%                                          %%
%% I ALTERED THE FILE usamts.tex            %%
%%                                          %%
%% If you have any questions or problems    %%
%% using this file, or with LaTeX in        %%
%% general, please go to the LaTeX          %%
%% forum on the Art Of Problem Solving      %%
%% web site, and post your problem.         %%
%%                                          %%
%%%%%%%%%%%%%%%%%%%%%%%%%%%%%%%%%%%%%%%%%%%%%%

%%%%%%%%%%%%%%%%%%%%%%%%%%%%%%%%%%%%%%%%%%%%
%% I ALTERED THE FOLLOWING LINES
\documentclass[12pt]{article}
\usepackage{amsmath,amssymb,amsthm}
\usepackage[pdftex]{graphicx}
\usepackage[inline]{asymptote}
\usepackage{fancyhdr}
\usepackage[english]{babel}
\pagestyle{fancy}
\usepackage[nottoc]{tocbibind}
\usepackage{hyperref}
\usepackage{setspace}
\usepackage{csquotes}
\usepackage[style=mla-new]{biblatex}
\include{usamts}

\renewcommand{\baselinestretch}{1.44}
%% I ALTERED THE ABOVE LINES
%%%%%%%%%%%%%%%%%%%%%%%%%%%%%%%%%%%%%%%%%%%%


%% If you would like to use Asymptote within this document (which is optional), 
%% you can find out how at the following URL:
%%
%%   http://www.artofproblemsolving.com/Wiki/index.php/Asymptote:_Advanced_Configuration
%%
%% As explained there, you will want to uncomment the line below.  But be
%% sure to check the website because there are several other steps that must 
%% be followed.
%% \usepackage{asymptote}


%% Enter your real name here
%% Example: \realname{David Patrick}
\realname{[REAL NAME]}

%% Enter your USAMTS username here
%% Example: \username{DPatrick}
%% IMPORTANT
%% If your username contains one of the following characters:
%%      # $ ~ _ ^ % { } &
%% then this character must be preceded by a backslash \
%% for example: if your user name is math_genius, then the line below should be
%% \username{math\_genius}

\username{[USERNAME]}

%% Enter your USAMTS ID# here
%% Example: \usamtsid{9999}
\usamtsid{[00000]}

%% Make sure that the year and round number are correct.
%% Year 27 is the academic year 2015-16
\usamtsyear{32}
\usamtsround{3}

%%
%% DO NOT ALTER THE FOLLOWING LINE
\begin{document}
%% DO NOT ALTER THE PREVIOUS LINE
%%

%%%%%%%%%%%%%%%%%%%%%%%%%%%%%%%%%%%%%%%%%%%%%%%%%%%%%%%%
%%
%% All solutions go in between the \begin{solution} and
%% the \end{solution} corresponding to the problem
%% number.  
%%
%% For example, suppose Problem 2 is 
%% "Solve for x: x + 3 = 5".
%% Your solution would look like this:
%%
%% \begin{solution}{2}
%% If $x+3=5$, then subtract 3 from both
%% sides of the equation to get $x=5-3=2$, so the 
%% solution is $x=2$.
%% \end{solution}
%%
%%%%%%%%%%%%%%%%%%%%%%%%%%%%%%%%%%%%%%%%%%%%%%%%%%%%%%%%

\begin{solution}{1} 
Hello! This is a USAMTS template. Please make a copy of this to put your solutions in. Diagram for P1 below.

% Credit to users cw357 and Constance-Variance on AoPS
\begin{asy}
size(10cm);
real s= 10.0;
int[][] x = {
{0,0,0,0,0},
{0,0,0,0,0},
{0,0,0,0,0},
{0,0,0,0,0},
{0,0,0,0,0}};
void square(int a, int b)
{
fill(s*(a,b)--s*(a+1,b)--s*(a+1,b+1)--s*(a,b+1)--cycle);
}
square(1,2);
square(1,3);
square(3,1);
square(3,2);
for(int i = 0; i < 6; ++i) {
draw(s*(i,0)--s*(i,5));
}
for(int i = 0; i < 6; ++i) {
draw(s*(0,i)--s*(5,i));
}
for(int k = 0; k<5; ++k){
for(int l = 0; l<5; ++l){
if(x[k][l]!=0){
label(scale(3.0)*string(x[k][l]),s*(l+0.5,-k+4.5));
}
}
}
void sudokuLabel(int p, int q, int r) {
label(string(r), s*(p, q) + (1, -1));
}
sudokuLabel(1, 1, 4);
sudokuLabel(2, 1, 1);
sudokuLabel(3, 1, 1);
sudokuLabel(4, 1, 3);
sudokuLabel(0, 3, 9);
sudokuLabel(2, 3, 9);
sudokuLabel(4, 3, 5);
sudokuLabel(0, 5, 3);
sudokuLabel(1, 5, 1);
sudokuLabel(2, 5, 3);
sudokuLabel(3, 5, 2);
\end{asy}

\end{solution}

\begin{solution}{2}
%% Solution to Problem 2 goes here
\end{solution}

\begin{solution}{3}
%% Solution to Problem 3 goes here
\end{solution}

\begin{solution}{4}

\end{solution}

\begin{solution}{5}
%% Solution to Problem 5 goes here
\end{solution}

%%
%% DO NOT ALTER THE FOLLOWING COMMAND
\end{document}
%% DO NOT ALTER THE PREVIOUS COMMAND
%%

%%%%%%%%%%%%%%%%%%%%%%%%%%%%%%%%%%%%%%%%%%%%%%
%%                                          %%
%% You may want to run LaTeX on this file   %%
%% to make sure that it compiles correctly  %%
%% before you submit it.                    %%
%%                                          %%
%% If you run LaTeX on this file before     %%
%% entering any solutions, the output file  %%
%% will be blank.  This is normal.          %%
%%                                          %%
%% Make sure that the file usamts.tex is    %%
%% in the same directory as this file.      %%
%% If the compiler hangs or asks for        %%
%% fancyhdr.sty when you compile, then go   %%
%% to the My Forms portion of My USAMTS     %%
%% at www.usamts.org to download the        %%
%% fancyhdr.sty file.                       %%
%%                                          %%
%% Questions?  Problems?  Go to the LaTeX   %%
%% forum on www.artofproblemsolving.com     %%
%%                                          %%
%%%%%%%%%%%%%%%%%%%%%%%%%%%%%%%%%%%%%%%%%%%%%%


\end{document}